\scanpage{53}

% archive sheet number: 52

\deleted{$r$ \inserted{\unsure{aux $\rho$}} indique que l'on utilise l'orientation définie par $r$ pour définir $\rho_s,\rho_{s'},\rho_{s''}$. \unsure{Une|Soit une} orientation \unsure{$\omega$} de $\Pi$, et une f.}

une suite (bien déterminée) $(s,s',s'')$ d'él. \unsure{distinct|distincts} de $S$, de sorte que, si $\omega$ est l'orientation de $\Pi$ \deleted{telle que la face} \inserted{qui \unsure{induit|induise} sur} la face $f=\{s,\unsure{s'|r'},s''\}$ l'orientation opposée de celle définie par l'ordre précédent, on ait

\[
\deleted{\unsure{\rho}}\Sigma=(\rho_s^\omega,\rho_{s'}^\omega,\rho_{s''}^\omega).
\]

[\uline{NB} la donnée de $s,s',s''$ équivaut à celle d'un repère, savoir $r=\{s,\{s\unsure{,|\mathord{\,}}s'\}\unsure{_s|{}},\{s,s',s''\}\}$].

\figure{\begin{tikzpicture}\draw (0,0) .. controls (1.3,-.65) and (2.2,-1.59) .. (3.75,-2.15);\end{tikzpicture}}

Cas de $\unsure{\mathfrak S_4|\mathcal S_4}$

\uline{Proposition} Soit $\Pi$ un octaèdre comb. \inserted{orienté}, et $H\inserted{=\mathrm{Aut}^{+}(\Pi)}$ le gpe de ses autom orientés, $R=\mathrm{Rep}\,(\Pi)$, $R^{+}=\mathrm{Rep}^{+}(\Pi)$

a) Pour tout $r\in R$, soit $\Sigma_r$ le syst. de Schwartz correspondant \inserted{de type $(\deleted{\unsure{4}}3,\deleted{\unsure{3}}4,2)$}. Alors $\Sigma_r=\Sigma_{r'}$ sss $r$ et $r'$ sont antipodiques, et l'ens des $\Sigma_r$ ($r\in R$) est égal à l'ens des $\Sigma_r$ ($r\in R^{+}$) et en corr. $1$--$1$ avec $R^{+}$ -- donc un torseur sous $H$.

\margin{$24$ d'où $6\times24=144$ \deleted{\unsure{NB}} syst. de Schwartz (en appliquant $\mathfrak S_3$)}

b) L'application $s\unsure{\longmapsto|\mathrel{\vdash\!\dashv}}\rho_s$ définit une bijection de $S$ avec $ {}_4H$ (ens des él. d'ordre $4$ de $H$)
